% Tons de cores utilizados
\definecolor{mpgreen}{RGB}{0,102,51} % Tom de verde ajustado para combinar com o logo
\definecolor{mpgreen2}{RGB}{0,90,60} % Tom de verde mais escuro para o gradiente
\definecolor{darkgray}{RGB}{51,51,51} % Cinza escuro para textos

\thispagestyle{empty}

% Início do desenho da capa
\begin{tikzpicture}[remember picture,overlay]
    % Barra vertical decorativa com gradiente
    \shade[left color=mpgreen,right color=white,opacity=0.2] 
        (current page.north west) rectangle ([xshift=8cm]current page.south west);
    
    % Elementos decorativos
    \draw[mpgreen,opacity=0.2,line width=2pt] ([xshift=8.2cm,yshift=-2cm]current page.north west) -- ([xshift=\paperwidth-2cm,yshift=-2cm]current page.north east);
    \draw[mpgreen,opacity=0.2,line width=2pt] ([xshift=8.2cm,yshift=-2.3cm]current page.north west) -- ([xshift=\paperwidth-2cm,yshift=-2.3cm]current page.north east);
\end{tikzpicture}

\begin{center}
    \color{darkgray}
    \small
    \textsc{Instituto Federal de Educação, Ciência e Tecnologia de Minas Gerais} \\[0.2cm] 
    \textsc{Campus Formiga} \\[0.2cm]
    \textsc{Programa de Pós-Graduação em Administração} \\[0.2cm]
    \textsc{Mestrado Profissional em Administração}
    
    \vspace{2cm}
    
    % Inserção do logo
    % \includegraphics[width=0.3\textwidth]{logo.png}
    
    \vspace{2cm}
    
    \begin{minipage}{0.8\textwidth}
        \centering
        {\LARGE \textbf{Relatório Integral de Autoavaliação}} \\[0.3cm] 
        {\Large \textbf{Comissão de Autovaliação}} \\[0.8cm]
        {\Large Quadriênio 2021 - 2024}
    \end{minipage}
    
    \vfill
    
    \begin{tikzpicture}[remember picture,overlay]
        % Elemento decorativo inferior
        \draw[mpgreen,opacity=0.2,line width=2pt] ([xshift=2cm,yshift=3cm]current page.south west) -- ([xshift=\paperwidth-2cm,yshift=3cm]current page.south east);
    \end{tikzpicture}
    
    \large
    Formiga, Minas Gerais\\[0.3cm]
    2025
\end{center}




% tons de cores
\definecolor{mpgreen}{RGB}{0,90,60} % Mantendo o mesmo tom de verde da capa
\definecolor{darkgray}{RGB}{51,51,51} % Mantendo o mesmo cinza escuro

% formatação do título de cada capítulo
\makeatletter
\def\@makechapterhead#1{%
  \vspace*{50\p@}%
  {\parindent \z@ \normalfont
    % Elemento decorativo superior
    \begin{tikzpicture}
      \draw[mpgreen!60!black,opacity=0.3,line width=4pt] 
        (0,0) -- (4,0);
    \end{tikzpicture}
    
    \vspace{1em}%
    {\color{darkgray}%
    \ifnum \c@secnumdepth >\m@ne
      \huge\bfseries \thechapter\space
    \fi
    \huge \bfseries #1\par\nobreak}
    
    \vspace{0.5em}%
    % Elemento decorativo inferior
    \begin{tikzpicture}
      \draw[mpgreen!60!black,opacity=0.3,line width=2pt] 
        (0,0) -- (2,0);
    \end{tikzpicture}
    \vskip 40\p@
  }}
\makeatother




\pagebreak

\thispagestyle{empty}

\begin{center}
    \vspace*{2cm}
    {\Large\textbf{Comissão de Autovaliação}}\\
    \vspace{0.3cm}
    {\Large\textbf{Programa de Pós-Graduação em
Administração do IFMG - Campus Formiga}}
    
    \vspace{0.5cm}
    {\small Instituída pela Portaria Nº 164 de 15 de Outubro de 2024}
    
    \vspace{2cm}
    
    \begin{minipage}{0.8\textwidth}
        \centering
        \textbf{Washington Santos da Silva}\\
        Presidente - Representante Docente
        
        \vspace{1cm}
        
        \textbf{Elaine Belo Veloso da Silva}\\
        Representante Técnico-Administrativo
        
        \vspace{1cm}
        
        \textbf{Lelis Pedro de Andrade}\\
        Representante Docente
        
        \vspace{1cm}
        
        \textbf{Igor Neves Nunes}\\
        Representante Discente
        
        \vspace{1cm}
        
        \textbf{Suelem Correia Garcia}\\
        Representante dos Egressos
        
        \vspace{1cm}
        
        \textbf{Carolina Greco}\\
        Representante da Comunidade Externa
    \end{minipage}
\end{center}


\newpage
% RESUMO
\thispagestyle{empty}

\newenvironment{meuresumo}{
  \clearpage
  \small
  \vspace{-1cm}
  \begin{center}
    \bfseries RESUMO
    \vspace{0.5em}
  \end{center}
  \begin{quote}
}{
  \end{quote}
  \vspace{-1.1em}
  \begin{center}
  \begin{minipage}{0.87\textwidth} 
  \textbf{Palavras-chave:} Comissão de Autovaliação. 
  Mestrado Profissional em Administração. Política de de Autoavaliação. 
  Relatório de Autoavaliação.
  \end{minipage}
  \end{center}
  \clearpage
}

\begin{meuresumo}
Este Relatório Integral de Autoavaliação do Programa de Pós-Graduação em 
Administração do IFMG – Campus Formiga, referente ao quadriênio 2021-2024, 
apresenta uma análise do desempenho do Mestrado Profissional em Administração 
com base nas avaliações e percepções dos diferentes segmentos acadêmicos. 
Elaborado em conformidade com a Política de Autoavaliação do Programa, 
que estabeleceu diretrizes para um processo avaliativo estruturado e 
contínuo. Os resultados evidenciam a resiliência e o compromisso do 
programa, iniciado em 2021 em um cenário desafiador marcado pela 
pandemia da COVID-19 e pela queda na demanda por cursos de mestrado e 
doutorado, com a formação de pesquisadores qualificados e com o 
atendimento das demandas do mercado e da sociedade. A metodologia adotada 
envolveu questionários com itens quantitativos e qualitativos, abordando 
dimensõescomo qualidade geral do programa, infraestrutura, gestão 
administrativa, internacionalização e atividades extracurriculares, entre 
outras. A análise destacou pontos fortes, como a qualificação do corpo 
docente, a adequação de diversos aspectos da matriz curricular e da 
infraestrutura disponível, entre outros, mas também identificou oportunidades de 
aprimoramento, incluindo a revisão curricular, a ampliação das atividades 
extracurriculares, o fortalecimento da internacionalização e melhorias 
na infraestrutura. Destaca-se que algunas demandas identificidas durante 
o processo avaliativo produziram ações concretas de melhoria do 
programa. O relatório buscou estabelecer a autoavaliação como uma 
ferramenta estratégica para a evolução contínuo do programa, 
propondo ações concretas para seu desenvolvimento e sustentabilidade. 
A implementação dessas ações será monitorada em ciclos avaliativos 
subsequentes, para que contribua efetivamente para a consolidação 
e evolução do curso. 
\end{meuresumo}



% Reinicia a numeração de páginas e define o estilo arábico
\clearpage
\pagenumbering{arabic}


